\section{Calibration horizons and variance-sensitive alternatives}
The calibration example is a stress test of a conditional certificate,
not a recommendation to execute 100,000 suites before changing a layout.
The main article's synthetic adoption table samples dependent prefixes
within each seed; its counts do not estimate the minimum required sample size.

For a frozen catalogue and $m\ge2$ IID epochs, a standard empirical-Bernstein
bound~\cite{maurerpontil}, applied to the negated and range-normalized
paired saving, replaces the Hoeffding radius by
\[
t^{\rm EB}_j=\sqrt{\frac{2\widehat v_j\log(2K/\delta)}m}
             +\frac{7R_j\log(2K/\delta)}{3(m-1)},
\]
where $\widehat v_j$ is the unbiased sample variance in squared bytes.
Using $\max\{A_j,\bar d_j-t^{\rm EB}_j-\rho R_j\}-I_j/h$ then gives
the same conditional TV/lifecycle conclusion. The finite-range term
remains even when the observed variance is zero. With $R_j=3584$ bytes,
$K=1$, $\delta=.05$ and $m=50$, that term alone is approximately
629.6 bytes, larger than the constructed 32-byte mean benefit. This
particular bound therefore does not make the example practical at 50 epochs;
it is an analytical comparison, not a new deployed or implemented certificate.
A tighter externally justified support range or a substantially larger
paired benefit can change the calculation. The frozen Hoeffding implementation
and all original uncertified outcomes remain unchanged.

A Wasserstein or moment-based allowance would additionally require a
declared metric or moment class, an appropriate regularity bound for the
packet-cost difference, and independently justified uncertainty for deployment
shift. Replacing the TV symbol alone supplies none of that information.
Variance sensitivity addresses calibration noise, not missing independence,
consumer utility, installation cost or an unsupported drift model.
